\documentclass[12pt,a4paper]{article}

\usepackage[T1]{fontenc}
\usepackage{lmodern}
\usepackage{amsmath}
\usepackage{graphicx}
\usepackage{geometry}
\usepackage{microtype}
\usepackage[hidelinks]{hyperref}
\usepackage{enumitem}

\geometry{margin=1in}
\setlength{\parindent}{0pt}
\setlength{\parskip}{0.7em}

\title{An Interactive Introduction to \LaTeX\\[0.3em]
\large Part 1: The Basics}
\author{Prepared by: Your Name\\
\small Based on the introductory material by Dr John D. Lees-Miller}
\date{\today}

\begin{document}

\maketitle
\tableofcontents
\newpage

\section{Why Use \LaTeX?}

\LaTeX\ is a document-preparation system designed to produce
well-formatted documents. It is especially useful for mathematics
and scientific writing. It was created by scientists for scientists,
has a large and active community, and can be extended through
thousands of packages. These packages support many kinds of work,
including papers, presentations, spreadsheets, and other technical
documents.

The main idea is to describe \emph{what the content is} rather than
manually controlling \emph{how every part looks}. The writer can
therefore focus on the content and allow \LaTeX\ to manage the
formatting.

\section{How \LaTeX\ Works}

A \LaTeX\ document is written as plain text containing commands that
describe the structure and meaning of the content. The \LaTeX\
program processes the source file and produces a formatted document,
usually as a PDF.

For example, the following source code emphasizes one word:

\begin{verbatim}
The rain in Spain falls \emph{mainly} on the plain.
\end{verbatim}

Its output is:

The rain in Spain falls \emph{mainly} on the plain.

Commands can also create lists, figures, equations, and many other
document elements.

\section{A Minimal \LaTeX\ Document}

A minimal \LaTeX\ document has the following form:

\begin{verbatim}
\documentclass{article}

\begin{document}

Hello World! % Your content goes here.

\end{document}
\end{verbatim}

Important ideas are:

\begin{itemize}
    \item A command normally begins with a backslash, such as
    \verb|\documentclass|.
    \item Every document starts with a \verb|\documentclass| command.
    \item The argument inside curly braces tells \LaTeX\ which kind of
    document to create. In this example, the class is \verb|article|.
    \item The main content is written between
    \verb|\begin{document}| and \verb|\end{document}|.
    \item A percent sign starts a comment. \LaTeX\ ignores the
    remainder of that source-code line.
\end{itemize}

\section{Writing and Compiling the Document}

An online editor such as Overleaf can be used to write and compile a
\LaTeX\ document. It compiles the source automatically and displays
the result. The examples in this document can be tested by typing them
into a \texttt{.tex} file and compiling it.

\section{Typesetting Ordinary Text}

For most ordinary text, the writer can type normally between
\verb|\begin{document}| and \verb|\end{document}|.

Words are separated by one or more spaces. However, several spaces in
the source are normally collapsed into a single space in the output.

Paragraphs are separated by one or more blank lines. For example, the
following source contains two paragraphs:

\begin{verbatim}
This is the first paragraph.

This is the second paragraph.
\end{verbatim}

A line break in the source does not necessarily create a line break in
the final PDF. For example,

\begin{verbatim}
The rain in Spain
falls mainly on the plain.
\end{verbatim}

is normally typeset as one continuous sentence.

\section{Quotation Marks and Special Characters}

\subsection{Quotation Marks}

Quotation marks require special source-code forms. A left single
quotation mark is written with a backtick, and a right single
quotation mark is written with an apostrophe:

\begin{verbatim}
`single quotation'
\end{verbatim}

Double quotation marks are written with two backticks on the left and
two apostrophes on the right:

\begin{verbatim}
``double quotation''
\end{verbatim}

The resulting output is `single quotation' and ``double quotation''.

\subsection{Special Characters}

Some characters have special meanings in \LaTeX. Common examples are:

\begin{center}
\begin{tabular}{ll}
\verb|%| & starts a comment,\\
\verb|#| & is used with macro parameters,\\
\verb|&| & marks an alignment point,\\
\verb|$| & begins or ends inline mathematics.
\end{tabular}
\end{center}

To display these characters as ordinary text, they are usually escaped
with a backslash:

\begin{verbatim}
\% \# \& \$
\end{verbatim}

The output is: \% \# \& \$.

\section{Basic Commands and Environments}

\subsection{Emphasis}

The command \verb|\emph| emphasizes text:

\begin{verbatim}
This word is \emph{important}.
\end{verbatim}

Output: This word is \emph{important}.

\subsection{A Bulleted List}

The \texttt{itemize} environment creates a bulleted list:

\begin{verbatim}
\begin{itemize}
    \item Tea
    \item Milk
    \item Biscuits
\end{itemize}
\end{verbatim}

The output is:

\begin{itemize}
    \item Tea
    \item Milk
    \item Biscuits
\end{itemize}

\subsection{A Figure}

A figure can be inserted with the \texttt{figure} environment and the
\verb|\includegraphics| command:

\begin{verbatim}
\begin{figure}
    \centering
    \includegraphics{gerbil}
\end{figure}
\end{verbatim}

The following compiling example displays \texttt{gerbil.png} when that
file is available; otherwise, it displays a placeholder:

\begin{figure}[htbp]
    \centering
    \IfFileExists{gerbil.png}
      {\includegraphics[width=0.45\textwidth]{gerbil.png}}
      {\fbox{\parbox[c][4cm][c]{0.45\textwidth}
      {\centering Place \texttt{gerbil.png} in the project folder.}}}
    \caption{An example of the \texttt{figure} environment.}
\end{figure}

\subsection{A Numbered Equation}

The \texttt{equation} environment displays and numbers an equation:

\begin{equation}
    \alpha+\beta+1
\end{equation}

\section{Handling Errors}

A compilation error occurs when \LaTeX\ cannot understand part of the
source. For example, misspelling \verb|\emph| as \verb|\meph| causes
an ``undefined control sequence'' error because \LaTeX\ does not know
the command \verb|\meph|.

Useful advice for handling errors is:

\begin{enumerate}
    \item Do not panic; errors are a normal part of writing source code.
    \item Fix errors as soon as they appear. The most recently added
    source is often the best place to begin checking.
    \item When several errors are reported, start with the first one.
    A later error may be only a consequence of an earlier mistake.
    \item Check the lines above the reported location because the real
    cause may appear earlier.
\end{enumerate}

\section{Typesetting Exercise 1: Text and Special Characters}

The following paragraph contains dollar signs, a percentage sign, and
quotation marks, so the special characters must be written correctly:

\begin{quote}
In March 2006, Congress raised that ceiling an additional
\$0.79 trillion to \$8.97 trillion, which is approximately
68\% of GDP. As of October 4, 2008, the
``Emergency Economic Stabilization Act of 2008''
raised the current debt ceiling to \$11.3 trillion.
\end{quote}

A suitable source form is:

\begin{verbatim}
In March 2006, Congress raised that ceiling an additional
\$0.79 trillion to \$8.97 trillion, which is approximately
68\% of GDP. As of October 4, 2008, the
``Emergency Economic Stabilization Act of 2008''
raised the current debt ceiling to \$11.3 trillion.
\end{verbatim}

\section{Inline Mathematics}

Dollar signs are used in pairs to mark mathematics inside a sentence.
For example:

\begin{verbatim}
Let $a$ and $b$ be distinct positive integers, and let
$c=a-b+1$.
\end{verbatim}

The output is: Let $a$ and $b$ be distinct positive integers, and let
$c=a-b+1$.

Mathematical expressions should be placed in math mode. Compare the
ordinary-text expression \texttt{c = a - b + 1} with the properly
typeset expression $c=a-b+1$.

\LaTeX\ manages mathematical spacing automatically. Therefore,

\begin{verbatim}
$y=mx+b$
\end{verbatim}

and

\begin{verbatim}
$y = m x + b$
\end{verbatim}

produce essentially the same mathematical spacing: $y=mx+b$.

Inline mathematics can also be written with the \texttt{math}
environment:

\begin{verbatim}
\begin{math}
y=mx+b
\end{math}
\end{verbatim}

\section{Superscripts, Subscripts, and Grouping}

A caret creates a superscript, and an underscore creates a subscript:

\begin{equation*}
    y=c_2x^2+c_1x+c_0.
\end{equation*}

Curly braces group several characters into one superscript or
subscript. For example, the correct Fibonacci relation is

\begin{equation*}
    F_n=F_{n-1}+F_{n-2}.
\end{equation*}

Without braces, \verb|F_n-1| means $F_n-1$, not $F_{n-1}$.

\section{Greek Letters and Common Mathematical Notation}

Greek letters are written with commands such as
\verb|\alpha|, \verb|\beta|, \verb|\mu|, \verb|\Omega|, and
\verb|\omega|.

For example,

\begin{equation*}
    \mu=Ae^{Q/RT},
\end{equation*}

where $\mu$, $A$, $Q$, $R$, and $T$ are mathematical symbols.

A summation is written with \verb|\sum|:

\begin{equation*}
    \Omega=\sum_{k=1}^{n}\omega_k.
\end{equation*}

Here, $k=1$ is the lower limit and $n$ is the upper limit.

\section{Displayed Equations}

A large mathematical expression can be displayed on its own line with
the \texttt{equation} environment. The roots of a quadratic equation
are given by

\begin{equation}
    x=\frac{-b\pm\sqrt{b^2-4ac}}{2a},
\end{equation}

where $a$, $b$, and $c$ are the coefficients of the quadratic
equation.

The commands used above include:

\begin{itemize}
    \item \verb|\frac{numerator}{denominator}| for a fraction,
    \item \verb|\pm| for the plus-or-minus symbol,
    \item \verb|\sqrt{...}| for a square root,
    \item \verb|^| for superscripts.
\end{itemize}

Blank lines should not be placed inside a mathematical environment.

\section{Understanding Environments}

An environment is a context created with matching \verb|\begin| and
\verb|\end| commands:

\begin{verbatim}
\begin{environment-name}
    Content
\end{environment-name}
\end{verbatim}

The same mathematical source can appear differently in different
contexts. For example, the expression
$\Omega=\sum_{k=1}^{n}\omega_k$ appears compactly in a sentence, while

\begin{equation}
    \Omega=\sum_{k=1}^{n}\omega_k
\end{equation}

appears larger in a displayed equation. The summation symbol and the
positions of its limits change according to the context.

\section{List Environments}

The \texttt{itemize} environment creates bullet points:

\begin{itemize}
    \item Biscuits
    \item Tea
\end{itemize}

The \texttt{enumerate} environment creates numbered items:

\begin{enumerate}
    \item Biscuits
    \item Tea
\end{enumerate}

Their source forms are:

\begin{verbatim}
\begin{itemize}
    \item Biscuits
    \item Tea
\end{itemize}

\begin{enumerate}
    \item Biscuits
    \item Tea
\end{enumerate}
\end{verbatim}

\section{Packages and the Preamble}

Many basic commands and environments are built into \LaTeX. Packages
are libraries that provide additional commands and environments.
Packages are loaded in the preamble, which is the part of the source
between \verb|\documentclass| and \verb|\begin{document}|.

For example, the \texttt{amsmath} package from the American
Mathematical Society is loaded as follows:

\begin{verbatim}
\documentclass{article}
\usepackage{amsmath}

\begin{document}

% Commands from amsmath can now be used here.

\end{document}
\end{verbatim}

This document also loads \texttt{graphicx} so that images can be
included.

\section{Unnumbered Equations with \texttt{amsmath}}

The starred \texttt{equation*} environment displays an equation
without assigning an equation number:

\begin{equation*}
    \Omega=\sum_{k=1}^{n}\omega_k.
\end{equation*}

Its source is:

\begin{verbatim}
\begin{equation*}
    \Omega=\sum_{k=1}^{n}\omega_k
\end{equation*}
\end{verbatim}

\section{Mathematical Operators}

In mathematics, adjacent letters are normally interpreted as separate
variables multiplied together. Therefore, named operators should use
appropriate operator commands.

A minimization expression should use \verb|\min|:

\begin{equation*}
    \min_{x,y}\left\{(1-x)^2+100(y-x^2)^2\right\}.
\end{equation*}

For operators without a predefined command, \verb|\operatorname| can
be used. For example,

\begin{equation*}
    \beta_i=
    \frac{\operatorname{Cov}(R_i,R_m)}
         {\operatorname{Var}(R_m)}.
\end{equation*}

In this expression, $\operatorname{Cov}$ denotes covariance and
$\operatorname{Var}$ denotes variance.

\section{Aligned Equations}

The \texttt{align*} environment aligns a sequence of equations. The
ampersand marks the alignment position, and a double backslash starts
a new line:

\begin{align*}
    (x+1)^3
        &= (x+1)(x+1)(x+1) \\
        &= (x+1)(x^2+2x+1) \\
        &= x^3+3x^2+3x+1.
\end{align*}

The corresponding source is:

\begin{verbatim}
\begin{align*}
(x+1)^3 &= (x+1)(x+1)(x+1) \\
        &= (x+1)(x^2+2x+1) \\
        &= x^3+3x^2+3x+1
\end{align*}
\end{verbatim}

Because the environment name contains a star, the equations are not
numbered.

\section{Typesetting Exercise 2: Probability and Mathematics}

Let $X_1,X_2,\ldots,X_n$ be a sequence of independent and identically
distributed random variables with
$\operatorname{E}[X_i]=\mu$ and
$\operatorname{Var}[X_i]=\sigma^2<\infty$, and let

\begin{equation*}
    S_n=\frac{1}{n}\sum_{i=1}^{n}X_i
\end{equation*}

denote their mean. Then, as $n$ approaches infinity, the random
variables

\begin{equation*}
    \sqrt{n}(S_n-\mu)
\end{equation*}

converge in distribution to a normal distribution
$\mathcal{N}(0,\sigma^2)$.

The command for the infinity symbol is \verb|\infty|. The source for
the main statement is:

\begin{verbatim}
Let $X_1,X_2,\ldots,X_n$ be independent and identically
distributed random variables with
$\operatorname{E}[X_i]=\mu$ and
$\operatorname{Var}[X_i]=\sigma^2<\infty$, and let

\begin{equation*}
S_n=\frac{1}{n}\sum_{i=1}^{n}X_i
\end{equation*}

denote their mean. Then, as $n$ approaches infinity,
$\sqrt{n}(S_n-\mu)$ converges in distribution to
$\mathcal{N}(0,\sigma^2)$.
\end{verbatim}

\section{Summary}

This document has demonstrated how to:

\begin{itemize}
    \item create a minimal \LaTeX\ article;
    \item write and format ordinary text;
    \item use comments, commands, arguments, and special characters;
    \item write correct quotation marks;
    \item emphasize text;
    \item create figures, bulleted lists, and numbered lists;
    \item recognize and respond to compilation errors;
    \item write inline and displayed mathematics;
    \item use superscripts, subscripts, grouping, Greek letters,
    fractions, square roots, sums, and infinity;
    \item work with environments;
    \item load packages in the preamble;
    \item create numbered and unnumbered equations;
    \item use mathematical operators;
    \item align multi-line equations.
\end{itemize}

A later stage of learning \LaTeX\ can introduce more structured
documents with sections, cross-references, figures, tables, and
bibliographies.

\end{document}
